Mathematical Specification of the Terminal Clearing Operator

\[
P_t^{(i)} \leftarrow \mathcal{A}\bigl(P_{t^-}^{(i)},\, D_t,\, \Omega_i\bigr)
\]

1. Definitions

\( P_{t^-}^{(i)} \in \mathbb{R}_+ \):  
  Pre-clearing price (or market value) of asset \(i\) in the old numeraire, observed at the instant immediately before the simultaneous-liquidation / PIK-saturation trigger.

\( D_t \in \mathbb{R}_+ \):  
  World deficit stock at time \(t\). Defined as the aggregate excess of nominal claims over realizable value across the entire visible claim structure:
  \[
  D_t = \sum_j C_j(t) - V_{\text{real}}(t)
  \]
  where \(C_j(t)\) are all outstanding nominal claims and \(V_{\text{real}}(t)\) is the contemporaneous valuation of the real-economy substrate. Once the trigger fires, \(D_t\) becomes the sole numeraire.

\( \Omega_i \in [0,1] \):  
  Residual claim priority / collateral quality of asset \(i\).  
  \(\Omega_i = 1\) corresponds to the highest-priority, best-collateralized residual claim;  
  \(\Omega_i = 0\) corresponds to pure residual / unsecured / fully subordinated claims.

2. Functional Form of the Clearing Algorithm \(\mathcal{A}\)

The operator \(\mathcal{A}\) performs a simultaneous, deficit-relative normalization of every asset. A canonical realization is the priority-weighted deficit clearing rule:

\[
\mathcal{A}\bigl(P_{t^-}^{(i)},\, D_t,\, \Omega_i\bigr)
=
\left(
\frac{P_{t^-}^{(i)}}{\sum_j P_{t^-}^{(j)} \cdot \Omega_j}
\right)
\cdot
\Omega_i
\cdot
D_t
\]

Interpretation

The term \(\dfrac{P_{t^-}^{(i)}}{\sum_j P_{t^-}^{(j)} \cdot \Omega_j}\) computes the relative weight of asset \(i\) in the pre-clearing universe, adjusted by priority.
Multiplication by \(\Omega_i\) further scales the claim by its residual seniority / collateral quality.
Multiplication by \(D_t\) re-denominates the result in units of the world deficit stock.

Thus every price is expressed as a pure fraction of the residual deficit, modulated by seniority.

3. Alternative Strict Seniority Form (Waterfall)

If absolute priority is enforced, the algorithm becomes a sequential waterfall:

\[
\begin{align*}
R_0 &= D_t \\
R_{k} &= R_{k-1} - \sum_{i \in \mathcal{S}k} P{t^-}^{(i)} \\
P_t^{(i)} &= 
\begin{cases}
P_{t^-}^{(i)} & \text{if } i \in \mathcal{S}k \text{ and } R{k-1} \ge \sum_{i \in \mathcal{S}k} P{t^-}^{(i)} \\
\dfrac{R_{k-1}}{\sum_{i \in \mathcal{S}k} P{t^-}^{(i)}} \cdot P_{t^-}^{(i)} & \text{if partial recovery} \\
0 & \text{if } R_{k-1} = 0
\end{cases}
\end{align*}
\]

where \(\mathcal{S}_k\) is the set of assets of seniority class \(k\), ordered from highest to lowest \(\Omega\).

4. Compact Operator Statement

\[
P_t^{(i)} 
= 
\mathcal{A}\bigl(P_{t^-}^{(i)},\, D_t,\, \Omega_i\bigr)
=
\frac{P_{t^-}^{(i)} \cdot \Omega_i}{\sum_j P_{t^-}^{(j)} \cdot \Omega_j} \cdot D_t
\]

This is the fully specified mathematical form of the terminal clearing algorithm. All assets are re-priced simultaneously in units of the world deficit stock \(D_t\), with relative recovery determined exclusively by residual claim priority \(\Omega_i\).
